% Multi-region transfer experiments (US-073, EPIC 11 Paper Track).
% \input-able section: no \documentclass / \begin{document} here.
% Prose in English (the manuscript is EN). Every number is rendered from the real
% EPIC 12 artefacts by scripts/build_us073_transfer_figures.py; nothing is hand-typed.
% Tables : paper/tables/us-073-transfer/{sen4agrinet_domain_gap,eurocropsml_kshot}.tex
% Figures: paper/figures/us-073-transfer/{kshot_curve,mexico_phenology}.{png,svg}
% Required packages in the host preamble: booktabs, graphicx, amsmath, hyperref, url.

\section{Multi-Region Transfer: Train in France, Extend Elsewhere}
\label{sec:multiregion}

A crop-mapping copilot trained only on French labels (PASTIS-R) is, by construction,
a French expert. This section turns that limitation into a measured capability: we
document the \emph{recipe} ``train in France $\rightarrow$ extend elsewhere'' along two
real transfer paths that were executed end-to-end and versioned with DVC, plus a
qualitative zero-shot demonstration over Mexico and an explicitly out-of-scope tropical
extension. The thesis we defend is deliberately conservative: the model \emph{receives
new classes from other datasets via few-shot adaptation}, and the Franco-Iberian domain
gap is \emph{catastrophic in zero-shot}. Consequently, the \textbf{$\Delta$mIoU between
zero-shot and few-shot is the scientific deliverable}, not a high absolute accuracy. This
is consistent with the limited spatial transferability of frozen Earth-observation
embeddings reported in ``Harvesting AlphaEarth''~\cite{harvesting2026alphaearth}: one must
budget for few-shot finetuning rather than promise zero-shot magic.

All embeddings are AlphaEarth Foundations~\cite{brown2025alphaearth} (Google DeepMind),
accessed through the Google Earth Engine asset \texttt{SATELLITE\_EMBEDDING/V1/ANNUAL},
data version~1.1, 64-dimensional, CC-BY-4.0. The frozen reasoner that turns model outputs
into explanations is Gemini~2.5~Pro~\cite{google2025gemini25} (GA, 1M-token context);
following the ``Be My Eyes'' pattern~\cite{huang2025bemyeyes}, our trained models are the
\emph{perceiver} and the LLM is a frozen \emph{reasoner} for communication and data
sovereignty, never a pixel classifier.

\subsection{Label-Space Harmonization (HCAT v3)}
\label{sec:multiregion:labelspace}

Transfer across datasets requires a common label space, otherwise the per-region heads
fight over incompatible gradients. We map the 18 agronomic PASTIS-R classes (the
contiguous \texttt{semantic18} space, ids $0\ldots17$) to HCAT~v3 leaf codes (10-digit,
prefix \texttt{33}~$=$~\texttt{crop\_type}) and collapse them, via the group level of the
HCAT hierarchy, to \textbf{11 canonical macro-classes} (10 crop groups plus a unified
\texttt{void}/background class), well inside the requested 10--20 range. Every HCAT leaf
code is verified to appear verbatim in the EuroCrops HCAT~v3 reference
(\texttt{data/reference/eurocrops\_hcat3.csv}); no code is invented. The mapping and its
collapse rule are documented in \texttt{docs/data/hcat\_crosswalk.md} and materialized as
\texttt{data/reference/hcat\_crosswalk.parquet}.

Three PASTIS classes have no 1:1 HCAT leaf and are flagged \texttt{approx} as explicit
agronomic decisions (not matching failures): \emph{Mixed cereal} maps to the L2 group node
\texttt{cereal}; \emph{Fruits/vegetables/flowers} maps to the dominant representative leaf
\texttt{fresh\_vegetables}; \emph{Leguminous fodder} maps to \texttt{legumes\_harvested\_green}.
The long tail is mitigated by the collapse: \emph{Meadow} ($\sim$45\% of crop parcels) is
isolated as \texttt{grassland}, and the eight sibling cereals (including the
near-inseparable soft and durum winter wheats) are fused into \texttt{cereals}, reducing
the inter-class variance that depressed the 18-class macro-F1. The cross-dataset
convention is uniform: \texttt{background}/\texttt{void} $\rightarrow$
\texttt{ignore\_index = 255} (excluded from the confusion matrix), while a class present in
one dataset but absent in another is treated as \texttt{null-class} (an unknown,
partial-label tag), never as a hard negative.

\subsection{Dense Path: Sen4AgriNet Catalonia (France $\rightarrow$ Spain)}
\label{sec:multiregion:dense}

The best individual French segmenter, TSViT~\cite{tarasiou2023tsvit} with the
phenological branch (TSViT-pheno, fold-4 mIoU~0.6253 / F1~0.7500 on
PASTIS-R~\cite{garnot2021pastis}), is the source model. We finetune it on a small
Sen4AgriNet~\cite{sykas2022sen4agrinet} subset over the Catalonia tile 31TCG
(Sentinel-2, years 2019/2020),
projecting both heads onto the shared HCAT macro space. The subset is intentionally tiny
(10 training patches / 90 tiles; 20 validation patches / 180 tiles), versioned in
DVC/GCS (\texttt{data/sen4agrinet.dvc}); we never download the full multi-terabyte
dataset.

Table~\ref{tab:sen4agrinet_domain_gap} reports the result. The zero-shot transfer yields
an mIoU of exactly~$0.0000$. This is real, not a bug: the French TSViT-pheno projected
$18\rightarrow$macro predicts the Catalonia validation tiles almost exclusively as
\texttt{vegetables} ($\sim$567k px) and \texttt{grassland} ($\sim$416k px), two classes
whose support is \emph{zero} in the Spanish ground truth (dominated instead by
\texttt{cereals}, \texttt{legumes\_fodder}, \texttt{vineyard} and
\texttt{oilseed\_industrial}). With an all-zero confusion diagonal, the mIoU is
exactly~0 -- a textbook catastrophic Franco-Iberian domain gap. We also found that the
Sen4AgriNet Catalonia labels carry S4A crop codes outside the PASTIS-18 nomenclature;
these are honestly routed to \texttt{ignore\_index = 255}, yet the remaining
\texttt{cereals}/\texttt{legumes}/\texttt{vineyard}/\texttt{oilseed} ground truth is
sufficient for a valid measurement. Finetuning the macro head on the 10 Spanish patches
recovers a few-shot mIoU of~$0.2468$ (F1-macro~$0.3005$, pixel-accuracy~$0.9179$), i.e.
$\Delta$mIoU~$= +0.2468$ over 40 epochs. This finetuned model is the
``Franco-Iberian'' segmenter that adds olive, vineyard, sorghum and durum-wheat behaviour
to the French source. The full run is recorded in
\texttt{reports/segmentation/sen4agrinet\_transfer\_result.json} and reproduced by the
executed notebook \texttt{notebooks/segmentation/5c\_transfer\_sen4agrinet.ipynb}.

\begin{table}[t]
\centering
\caption{Dense France $\rightarrow$ Catalonia transfer (TSViT-pheno finetuned on the Sen4AgriNet 31TCG subset, HCAT macro label-space). Trained on 10 ES patches, evaluated on 20. The catastrophic zero-shot mIoU and the recovered few-shot mIoU yield $\Delta$mIoU $= +0.2468$ -- the measured Franco-Iberian domain gap is the deliverable, not a high accuracy. Source: \texttt{reports/segmentation/sen4agrinet\_transfer\_result.json}. Sen4AgriNet (Sykas et al. 2022, CC-BY-SA-4.0).}
\label{tab:sen4agrinet_domain_gap}
\begin{tabular}{lrrr}
\toprule
Protocol & mIoU & F1-macro & Pixel-acc. \\
\midrule
Zero-shot (FR 18 $\rightarrow$ macro) & 0.0000 & 0.0000 & 0.0000 \\
Few-shot (10 ES patches) & 0.2468 & 0.3005 & 0.9179 \\
\midrule
$\Delta$ (few-shot $-$ zero-shot) & $+0.2468$ & $+0.3005$ & $+0.9179$ \\
\bottomrule
\end{tabular}
\end{table}


\subsection{Tabular Few-Shot Path: EuroCropsML (Transnational)}
\label{sec:multiregion:fewshot}

The dense path measures a within-Europe domain gap pixel-wise; the tabular path measures
how many local labels close a \emph{transnational} gap. Two factual clarifications are
mandatory. First, \textbf{France is not part of EuroCropsML}: the benchmark~\cite{reuss2025eurocropsml} covers Estonia (EE), Latvia (LV) and Portugal (PT),
so the real protocol is LV$+$PT~$\rightarrow$~EE and LV~$\rightarrow$~EE (Reuter et al.
2025, Table~II), not ``France~$\rightarrow$~Estonia''.
Second, EuroCropsML carries \textbf{no AlphaEarth embedding}: each datapoint is an annual
Sentinel-2 L1C median time series (13 bands, B10 excluded) per parcel. We therefore reuse
the \emph{recipe} of our AlphaEarth tabular baseline (an XGBoost \texttt{multi:softprob}
over a fixed-width per-parcel vector) but build that vector from the parcel's own
Sentinel-2 series; the AlphaEarth-via-GEE variant for these countries stays future work.
Labels are aligned to the same HCAT macro space (the \texttt{EC\_hcat\_c} code is an HCAT
leaf code), yielding 8 macro-classes on this subset, versioned as
\texttt{data/transfer/eurocropsml.dvc}.

Table~\ref{tab:eurocropsml_kshot} and Figure~\ref{fig:eurocropsml_kshot} give the
F1-macro-versus-$k$ curve (mean~$\pm$~std over three seeds) for $k \in
\{1,5,10,20,100,200,500\}$ shots per target class, in three scenarios: source-pre-trained
LV$+$PT~$\rightarrow$~EE, source-pre-trained LV~$\rightarrow$~EE, and a no-pre-train
reference. The reading is the intended one for a domain gap: at $k{=}1$ the no-pre-train
model is near-random ($\sim$0.015 F1-macro) while source pre-training already reaches
$\sim$0.32; as $k$ grows the curves converge, confirming that a handful of local labels --
not a magical zero-shot -- is what closes the gap. The curve is evidence of the measured
transfer cost, never a zero-shot accuracy claim, and is computed strictly over the real
EuroCropsML splits.

\begin{table}[t]
\centering
\caption{EuroCropsML transnational few-shot transfer curve (XGBoost recipe on the parcel Sentinel-2 series, HCAT macro label-space, 8 classes). F1-macro mean $\pm$ std over 3 seeds per k-shot. France is NOT in EuroCropsML; the protocol is LV[+PT] $\rightarrow$ EE (Reuss et al. 2025, Table II), not France $\rightarrow$ Estonia. The no-pre-train column quantifies how source pre-training closes the gap at low k. EuroCropsML (Reuss et al. 2025, CC-BY-SA-4.0).}
\label{tab:eurocropsml_kshot}
\begin{tabular}{rccc}
\toprule
$k$ & LV+PT $\rightarrow$ EE (pre-train) & LV $\rightarrow$ EE (pre-train) & LV $\rightarrow$ EE (no pre-train) \\
\midrule
1 & $0.315 \pm 0.005$ & $0.321 \pm 0.003$ & $0.015 \pm 0.015$ \\
5 & $0.321 \pm 0.008$ & $0.349 \pm 0.022$ & $0.198 \pm 0.020$ \\
10 & $0.338 \pm 0.004$ & $0.354 \pm 0.020$ & $0.230 \pm 0.005$ \\
20 & $0.350 \pm 0.006$ & $0.382 \pm 0.015$ & $0.273 \pm 0.017$ \\
100 & $0.380 \pm 0.006$ & $0.444 \pm 0.006$ & $0.380 \pm 0.013$ \\
200 & $0.406 \pm 0.007$ & $0.464 \pm 0.001$ & $0.426 \pm 0.010$ \\
500 & $0.420 \pm 0.007$ & $0.482 \pm 0.012$ & $0.472 \pm 0.008$ \\
\bottomrule
\end{tabular}
\end{table}


\begin{figure}[t]
  \centering
  \includegraphics[width=0.9\linewidth]{figures/us-073-transfer/kshot_curve.png}
  \caption{EuroCropsML transnational few-shot curve (LV$[+$PT$]\rightarrow$EE). F1-macro
  on the Estonian query set versus $k$ shots per class, mean over three seeds with std
  error bars. Source pre-training dominates at low $k$ and all curves converge as local
  supervision grows: the gap is closed by few-shot data, not by zero-shot transfer.
  EuroCropsML (Reuter et al. 2025, CC-BY-SA-4.0).}
  \label{fig:eurocropsml_kshot}
\end{figure}

\subsection{Qualitative Zero-Shot Demonstration: Mexico}
\label{sec:multiregion:mexico}

To show the same pipeline replicates outside Europe, we apply it to two real Mexican
producing areas with \emph{no} curated ground truth: avocado around Uruapan/Tancitaro
(Michoacan) and guava around Calvillo (Aguascalientes). For each area-of-interest we pull
the real 64-dimensional AlphaEarth zonal embedding and the real Sentinel-2 NDVI time
series from Google Earth Engine, and align the phenological narrative to the
\texttt{PERMANENT\_WOODY} family of the HCAT macro space as a qualitative analogy.
Figure~\ref{fig:mexico_phenology} shows the two NDVI curves: both crops are perennial
woody (orchard-like), with a relatively high and stable signal rather than the
sow--harvest cycle of an annual cereal, and the guava AOI shows the expected late-season
modulation from irrigation and pruning. This demonstration is purely methodological:
because there is no ground truth, \textbf{we never report an F1, accuracy or classifier
over Mexico} -- a meta-test in the codebase enforces that no classification metric is
imported in this path. It illustrates the mechanism of ``receiving new classes'', not a
performance claim. AlphaEarth (Satellite Embedding V1 Annual, CC-BY-4.0);
contains modified Copernicus Sentinel data.

\begin{figure}[t]
  \centering
  \includegraphics[width=0.9\linewidth]{figures/us-073-transfer/mexico_phenology.png}
  \caption{Qualitative Mexico demonstration: real Sentinel-2 NDVI phenology of avocado
  (Uruapan, Michoacan) and guava (Calvillo, Aguascalientes), both perennial-woody crops.
  No ground truth exists for these AOIs, so no F1/accuracy is computed or implied. The
  figure illustrates the methodology of extending to a new region, not a metric.}
  \label{fig:mexico_phenology}
\end{figure}

\subsection{Future / Scale: Tropical Crops}
\label{sec:multiregion:future}

The recipe scales, in principle, to tropical agriculture (rice, soybean, sugarcane,
cotton) through WorldCereal RDM and the Harmonized Global Crops dataset. We declare this
\textbf{explicitly out of scope and future} for the present manuscript: it was not
demonstrated within this work's window, and we make no metric claim about it. The license
attributions are nonetheless pre-registered (WorldCereal RDM, CC-BY per collection;
Harmonized Global Crops, CC-BY-SA-4.0) so that the legal compliance of any future ingestion
is already resolved.

\subsection{Caveat: Limited Spatial Transferability}
\label{sec:multiregion:caveat}

Both real paths corroborate the central caveat of ``Harvesting
AlphaEarth''~\cite{harvesting2026alphaearth}: frozen global embeddings transfer poorly
across distant agronomic regions without local adaptation. The dense France~$\rightarrow$~Catalonia zero-shot
mIoU of~$0.0000$ is the strongest possible empirical statement of this gap, and the
EuroCropsML curve quantifies the few-shot budget needed to bridge it. We therefore do
\emph{not} claim that zero-shot works outside France, nor that the LLM layer ``beats''
Gemini at classification (its role is communication, explanation and on-premise data
sovereignty via Qwen3.5-35B-A3B~\cite{yang2025qwen3} served with vLLM GPTQ-Int4 on a
single GPU), nor any
$\geq$0.80 F1 in Mexico. The deliverable of this section is the \emph{measured}
$\Delta$mIoU and the few-shot curve -- honest evidence of how the copilot becomes a
multi-region expert.

% Share-alike note: figures and reports that combine derivatives of Sen4AgriNet,
% EuroCropsML and Harmonized Global Crops inherit the CC-BY-SA-4.0 clause and must be
% redistributed under a compatible license (see docs/licenses/DATA_LICENSE.md).
